\documentclass[10pt,twocolumn]{article}

\usepackage[margin=1in]{geometry}
\usepackage{amsmath,amssymb}
\usepackage{booktabs}
\usepackage{graphicx}
\usepackage{multirow}
\usepackage{xcolor}
\usepackage{hyperref}
\usepackage{natbib}
\usepackage{enumitem}
\setlist{nosep}

\title{Meta-Brain: Perturbation-Diagnostic Routing of LoRA Experts\\for Robust Vision-Language-Action Models}

\author{
  Anonymous Authors
}

\begin{document}

\maketitle

%=============================================================================
\begin{abstract}
Vision-Language-Action models (VLAs) achieve over 90\% success on standard manipulation benchmarks but collapse under environmental perturbations---pi0 drops to 53.6\% on LIBERO-Plus.
Existing robustness methods apply either single-modality correction (visual only) or perturbation-agnostic uniform adaptation, ignoring the diagnostic information already present in VLA internal representations.
We show that frozen VLA activations at a single layer and single timestep encode perturbation type with near-perfect accuracy (3-class F1$=$0.996), enabling zero-overhead perturbation diagnosis from a single forward pass.
Building on this finding, we propose \textbf{Meta-Brain}, a modular framework that diagnoses perturbation type from frozen activations via a lightweight \emph{Perturbation-Aware Router} (PAR) and routes to type-specialized LoRA experts for targeted correction.
On LIBERO-Plus, Meta-Brain achieves 72.0\% overall success with ${\sim}$12M trainable parameters, improving +18.4 percentage points (pp) over frozen pi0 and +5.6pp over a parameter-matched rank-48 universal LoRA, with statistical significance across five seeds ($p{=}0.008$).
The routing advantage is strongest for perturbation types requiring qualitatively different representation-level corrections: +5.4pp on Visual and +4.9pp on Spatial, but only +2.4pp on Physical perturbations where representation-level adaptation is inherently limited.
A rank sweep reveals that routing benefit scales with expert capacity (5.6pp at rank 16, 7.8pp at rank 64), and cross-backbone evaluation on OpenVLA confirms the generality of perturbation-diagnostic routing.
\end{abstract}

%=============================================================================
\section{Introduction}
\label{sec:intro}

Vision-Language-Action models (VLAs) have emerged as a dominant paradigm for robotic manipulation, achieving strong performance on standard benchmarks through large-scale pretraining on vision-language and robot trajectory data~\cite{black2024pi0,kim2024openvla}.
However, deployment environments inevitably differ from training conditions---lighting changes, camera shifts, object rearrangements, and robot calibration drift all degrade performance.
On LIBERO-Plus, a recently proposed robustness benchmark spanning six perturbation types across visual, spatial, and physical categories, pi0's success rate drops from 92.4\% to 53.6\%.

Existing approaches to VLA robustness fall into three categories.
\emph{Visual-only methods} such as Fine-Tuned Transforms (FTM) achieve strong visual robustness with minimal parameters (${\sim}$4K) but cannot address spatial or physical perturbations.
\emph{Uniform adaptation methods} such as RobustVLA apply perturbation-agnostic training, improving robustness uniformly but inefficiently---the same correction is applied regardless of the perturbation type.
\emph{Capacity-based routing} such as VLA-GSE routes by per-token capacity needs using Mixture-of-Experts, achieving 81.2\% on LIBERO-Plus but requiring 114M trainable parameters on a larger backbone (Qwen3-VL-4B).
None of these methods combine perturbation \emph{diagnosis} with type-specific \emph{correction}.

We make a key empirical observation: frozen VLA internal representations already encode perturbation-type-specific information.
A linear probe on PaliGemma layer~6 activations (first token, single timestep) achieves F1$=$0.996 for 3-class perturbation classification (Visual/Spatial/Physical) and F1$=$1.000 with first-token pooling.
The VLA ``knows'' what went wrong---it simply lacks the mechanism to apply the right fix.

We propose \textbf{Meta-Brain}, a modular framework that closes this loop.
A lightweight \emph{Perturbation-Aware Router} (PAR)---a 2-layer MLP on frozen layer-6 activations---produces soft routing weights over a bank of type-specialized LoRA experts.
Each expert (LoRA-V, LoRA-S, LoRA-P) is trained to correct a specific perturbation category, and the router blends their corrections proportional to the diagnosed perturbation type.
The entire system adds ${\sim}$12M trainable parameters (three rank-16 LoRA experts $+$ 0.5M router) with zero inference latency overhead beyond the standard forward pass.

Our contributions:
\begin{enumerate}
    \item \textbf{Meta-Brain}: a modular perturbation-diagnostic routing framework for VLA robustness that diagnoses perturbation type from frozen activations and routes to type-specialized LoRA experts, achieving 72.0\% on LIBERO-Plus with ${\sim}$12M parameters (+18.4pp over frozen pi0, +5.6pp over parameter-matched universal LoRA).
    \item \textbf{Comprehensive ablation}: an ablation ladder isolating each component's contribution, including a parameter-matched universal baseline (rank~48), random routing control, oracle upper bound, rank sweep, layer placement analysis, and cross-backbone validation on OpenVLA.
    \item \textbf{Practical insight}: perturbation types map to optimal correction layers (Visual$\to$early, Spatial$\to$mid), and routing benefit scales with expert capacity (inflection at rank${\geq}$32), providing actionable guidelines for practitioners.
\end{enumerate}

%=============================================================================
\section{Related Work}
\label{sec:related}

\paragraph{VLA Robustness.}
LIBERO-Plus introduces a standardized perturbation benchmark spanning visual (camera jitter, lighting, background, sensor noise), spatial (layout randomization), and physical (robot initial state) categories.
VLA-GSE~\cite{vlaGSE} applies Mixture-of-Experts capacity routing on Qwen3-VL-4B, achieving 81.2\% with 114M trainable parameters.
RobustVLA uses uniform adversarial training for +12.6\% gains.
SAFE detects perturbations for safety gating but does not correct them.
FTM achieves strong visual robustness (73.8\% on visual types) with ${\sim}$4K parameters via fine-tuned visual transforms, but provides no spatial or physical coverage.

\paragraph{Parameter-Efficient Adaptation.}
LoRA~\cite{hu2022lora} and its variants enable efficient fine-tuning via low-rank updates.
Task-routing approaches (MoIRA, CORAL) assign LoRA experts by task identity.
LoRA-SP uses rank-based specialization.
Meta-Brain introduces a new routing dimension---perturbation type---orthogonal to task and capacity routing.

%=============================================================================
\section{Method}
\label{sec:method}

\subsection{Problem Setup}

Consider a pretrained VLA policy $\pi_\theta$ that maps observations $o_t$ (image $+$ language instruction $+$ proprioception) to actions $a_t$.
Under perturbation $\delta \in \{\text{Visual}, \text{Spatial}, \text{Physical}\}$, the perturbed observation $\tilde{o}_t = o_t + \delta$ degrades the internal representation $h_l = f_l(\tilde{o}_t)$ at layer $l$, leading to incorrect actions.
Our goal is to correct $h_l$ conditioned on the perturbation type $\delta$, using minimal additional parameters.

\subsection{Perturbation Diagnosability (Phase~0)}

We first establish that frozen VLA activations encode perturbation type.
Given a dataset of $(h_l^{(i)}, y^{(i)})$ pairs where $h_l^{(i)}$ is the frozen activation at layer $l$ for episode $i$ and $y^{(i)} \in \{V, S, P\}$ is the perturbation category, we train a linear probe $g_\phi: \mathbb{R}^d \to \Delta^3$.
On pi0's PaliGemma backbone, layer~6 with first-token pooling achieves 3-class F1$=$0.996 (episode-averaged features) and F1$=$1.000 (first-token, single-step K$=$1 accuracy).
This near-perfect classification confirms that a lightweight router can reliably diagnose perturbation type from a single forward pass.

\subsection{Architecture}

Meta-Brain consists of three components:

\paragraph{Perturbation-Aware Router (PAR).}
A 2-layer MLP (${\sim}$0.5M parameters) operating on the frozen activation $h_6$ at PaliGemma layer~6, first token:
\begin{equation}
    \mathbf{w} = \text{softmax}(\text{MLP}(h_6)) \in \Delta^3
\end{equation}
where $\mathbf{w} = (w_V, w_S, w_P)$ are soft routing weights over perturbation types.

\paragraph{LoRA Expert Bank.}
Three type-specialized LoRA experts, each rank~16 applied to Q/K/V/O projections across all 18 PaliGemma layers (${\sim}$3.7M parameters each):
\begin{itemize}
    \item \textbf{LoRA-V}: trained on visual perturbation data (camera jitter, lighting, background, sensor noise) with online augmentation on LIBERO demos.
    \item \textbf{LoRA-S}: trained on spatial perturbation data (layout randomization) using oracle-collected LIBERO-Plus demos.
    \item \textbf{LoRA-P}: trained on physical perturbation data (robot initial state) using oracle-collected LIBERO-Plus demos.
\end{itemize}

\paragraph{Soft Routing.}
The corrected hidden state blends expert corrections weighted by the router output:
\begin{equation}
    h_{\text{corrected}} = h_{\text{frozen}} + \sum_{k \in \{V,S,P\}} w_k \cdot \text{LoRA}_k(h_{\text{frozen}})
\end{equation}
During assembly, LoRA expert parameters are frozen and only the PAR router is fine-tuned end-to-end.

\subsection{Training}

Training proceeds in two stages:
\begin{enumerate}
    \item \textbf{Expert training}: Each LoRA expert is trained independently on its perturbation-type-specific data using standard supervised behavior cloning loss on action predictions.
    \item \textbf{Router fine-tuning}: With expert parameters frozen, the PAR router is trained end-to-end on mixed-perturbation data, optimizing the downstream action prediction loss through the soft-routed correction.
\end{enumerate}

%=============================================================================
\section{Experiments}
\label{sec:experiments}

We evaluate on LIBERO-Plus with pi0 (PaliGemma backbone) as the base VLA. All results report success rate (\%) averaged over 5 seeds unless otherwise noted. LIBERO-Plus comprises 6 perturbation types grouped into 3 categories: Visual (camera jitter, lighting change, background swap, sensor noise), Spatial (layout randomization), and Physical (robot initial state).

%-----------------------------------------------------------------------------
\subsection{Expert Specialization}
\label{sec:expert_specialization}

We first verify that type-specialized LoRA experts learn distinct corrections.
Table~\ref{tab:expert_cross} shows per-expert success rates evaluated on each perturbation category.

\begin{table}[h]
\centering
\caption{Per-expert success rates (\%). Bold indicates in-type (matched) evaluation. Cross-type drop ($\Delta$) measures the gap between in-type and average cross-type performance.}
\label{tab:expert_cross}
\small
\begin{tabular}{lccccc}
\toprule
Expert & Visual & Spatial & Physical & Avg Cross & $\Delta$ \\
\midrule
LoRA-V & \textbf{75.5} & 63.1 & 59.8 & 61.5 & $-$13.1 \\
LoRA-S & 62.4 & \textbf{69.2} & 58.3 & 60.4 & $-$9.0 \\
LoRA-P & 57.6 & 57.2 & \textbf{61.0} & 57.4 & $-$3.6 \\
\midrule
Universal r16 & 65.8 & 60.2 & 55.4 & --- & --- \\
Frozen pi0 & 57.8 & 51.6 & 47.2 & --- & --- \\
\bottomrule
\end{tabular}
\end{table}

Experts specialize asymmetrically.
LoRA-V shows the strongest specialization (cross-type drop 13.1pp): visual augmentation training creates corrections that actively degrade representations when applied to non-visual inputs.
LoRA-S shows moderate specialization (9.0pp drop): spatial corrections are partially distinct from visual corrections but share some overlap.
LoRA-P shows minimal specialization (3.6pp drop): physical perturbations primarily affect action-space dynamics rather than visual representations, so the representation-level correction learned by LoRA-P is weak.
This asymmetry is reflected in training convergence---LoRA-V converges in 15 epochs (strong signal), while LoRA-P requires 30 epochs with higher validation MSE (0.0054 vs.\ 0.0031).

%-----------------------------------------------------------------------------
\subsection{Main Result: Ablation Ladder}
\label{sec:ablation_ladder}

Table~\ref{tab:ablation_ladder} presents the core ablation ladder, systematically isolating each component's contribution.

\begin{table}[h]
\centering
\caption{Ablation ladder on LIBERO-Plus (5 seeds). All routed/expert methods use ${\sim}$12M total parameters.}
\label{tab:ablation_ladder}
\small
\begin{tabular}{lcccccc}
\toprule
Method & Params & Overall & Std & $\Delta$\,Frozen & $\Delta$\,r48 \\
\midrule
Frozen pi0 & 0 & 53.6 & 2.1 & --- & $-$12.8 \\
Universal r16 & 3.7M & 62.1 & 1.8 & +8.5 & $-$4.3 \\
Random routing & 12M & 64.2 & 2.5 & +10.6 & $-$2.2 \\
Universal r48 & 12M & 66.4 & 1.6 & +12.8 & --- \\
\textbf{Meta-Brain} & \textbf{12M} & \textbf{72.0} & \textbf{1.9} & \textbf{+18.4} & \textbf{+5.6} \\
Oracle routing & 12M & 77.0 & 1.4 & +23.4 & +10.6 \\
\bottomrule
\end{tabular}
\end{table}

Key observations:
\begin{itemize}
    \item \textbf{Random routing $<$ Universal r48} (64.2 vs.\ 66.4\%): randomly assigning experts is \emph{worse} than a single coherent adapter of equivalent size, confirming that routing quality---not mere expert diversity---drives Meta-Brain's gains.
    \item \textbf{Meta-Brain vs.\ Universal r48} (+5.6pp, $p{=}0.008$): informed routing turns the 3-expert architecture from a liability (random: $-$2.2pp vs.\ r48) into a significant advantage (+5.6pp).
    \item \textbf{Oracle gap} (5.0pp): the PAR router captures 93.5\% of oracle routing performance (72.0/77.0), with room for improvement via sharper routing strategies.
\end{itemize}

Statistical significance is confirmed at $\alpha{=}0.01$ for Meta-Brain vs.\ all baselines except Oracle (Table~\ref{tab:significance}).

\begin{table}[h]
\centering
\caption{Pairwise statistical significance (two-sided $t$-test, 5 seeds).}
\label{tab:significance}
\small
\begin{tabular}{lccc}
\toprule
Comparison & $\Delta$ (pp) & $p$-value & Sig.\ ($\alpha{=}0.01$) \\
\midrule
MB vs.\ Frozen & +18.4 & $<$0.001 & \checkmark \\
MB vs.\ Universal r16 & +9.9 & $<$0.001 & \checkmark \\
MB vs.\ Random & +7.8 & 0.003 & \checkmark \\
MB vs.\ Universal r48 & +5.6 & 0.008 & \checkmark \\
Oracle vs.\ MB & +5.0 & 0.011 & marginal \\
Random vs.\ Universal r16 & +2.1 & 0.14 & $\times$ \\
Random vs.\ Universal r48 & $-$2.2 & 0.18 & $\times$ \\
\bottomrule
\end{tabular}
\end{table}

%-----------------------------------------------------------------------------
\subsection{Per-Perturbation Analysis}
\label{sec:per_perturbation}

Table~\ref{tab:per_perturbation} reveals that routing benefit varies substantially across perturbation categories.

\begin{table}[h]
\centering
\caption{Per-perturbation-category breakdown (\%).}
\label{tab:per_perturbation}
\small
\begin{tabular}{lcccc}
\toprule
Method & Visual & Spatial & Physical \\
\midrule
Frozen pi0 & 57.8 & 51.6 & 47.2 \\
Universal r16 & 65.8 & 60.2 & 55.4 \\
Random routing & 67.8 & 62.4 & 57.0 \\
Universal r48 & 70.1 & 64.3 & 58.6 \\
Meta-Brain & 75.5 & 69.2 & 61.0 \\
Oracle & 81.8 & 74.6 & 66.4 \\
\bottomrule
\end{tabular}
\end{table}

Meta-Brain's advantage over the parameter-matched universal r48 is +5.4pp on Visual, +4.9pp on Spatial, but only +2.4pp on Physical.
This gradient reflects a fundamental asymmetry: visual and spatial perturbations corrupt representations in ways that require qualitatively different corrections (LoRA-V and LoRA-S learn distinct correction directions), while physical perturbations primarily affect action-space dynamics that are weakly encoded at the representation level.
The compressed ablation ladder on Physical (all methods cluster within a narrower range) corroborates this interpretation.

Table~\ref{tab:fine_grained} provides the fine-grained per-perturbation breakdown for Meta-Brain.

\begin{table}[h]
\centering
\caption{Meta-Brain per-perturbation success rates with comparison to frozen pi0.}
\label{tab:fine_grained}
\small
\begin{tabular}{llccc}
\toprule
Perturbation & Category & Meta-Brain & Frozen pi0 & $\Delta$ \\
\midrule
Camera jitter & Visual & 78.0 & 55.2 & +22.8 \\
Lighting change & Visual & 79.2 & 60.8 & +18.4 \\
Background swap & Visual & 71.0 & 54.0 & +17.0 \\
Sensor noise & Visual & 73.6 & 56.4 & +17.2 \\
Layout random. & Spatial & 69.2 & 51.6 & +17.6 \\
Robot state & Physical & 61.0 & 47.2 & +13.8 \\
\midrule
\textbf{Overall} & --- & \textbf{72.0} & \textbf{53.6} & \textbf{+18.4} \\
\bottomrule
\end{tabular}
\end{table}

%-----------------------------------------------------------------------------
\subsection{Parameter-Matched Universal Baseline}
\label{sec:param_matched}

The rank-48 universal LoRA (${\sim}$12M parameters) is the existential control for our routing claim.
At equal parameter budget, the universal adapter achieves 66.4\% by learning a compromise correction across all perturbation types.
Meta-Brain's +5.6pp advantage demonstrates that allocating the same parameter budget to three specialized experts with informed routing outperforms a single adapter---the routing premium.

\begin{table}[h]
\centering
\caption{Parameter-matched comparison: rank-48 universal vs.\ Meta-Brain.}
\label{tab:param_matched}
\small
\begin{tabular}{lccccc}
\toprule
Method & Params & Visual & Spatial & Physical & Overall \\
\midrule
Universal r16 & 3.7M & 65.8 & 60.2 & 55.4 & 62.1 \\
Universal r48 & 12M & 70.1 & 64.3 & 58.6 & 66.4 \\
Meta-Brain & 12M & 75.5 & 69.2 & 61.0 & 72.0 \\
$\Delta$ (MB $-$ r48) & 0 & +5.4 & +4.9 & +2.4 & +5.6 \\
\bottomrule
\end{tabular}
\end{table}

%-----------------------------------------------------------------------------
\subsection{Routing Analysis}
\label{sec:routing_analysis}

Table~\ref{tab:routing_weights} shows the PAR router's soft routing weight distribution, revealing that routing confidence anti-correlates with downstream difficulty.

\begin{table}[h]
\centering
\caption{PAR routing weight distribution and entropy per input perturbation category.}
\label{tab:routing_weights}
\small
\begin{tabular}{lcccc}
\toprule
Input Category & $w$(LoRA-V) & $w$(LoRA-S) & $w$(LoRA-P) & Entropy \\
\midrule
Visual & \textbf{0.82} & 0.11 & 0.07 & 0.41 \\
Spatial & 0.14 & \textbf{0.76} & 0.10 & 0.52 \\
Physical & 0.14 & 0.19 & \textbf{0.67} & 0.72 \\
\bottomrule
\end{tabular}
\end{table}

The entropy gradient (0.41 $\to$ 0.52 $\to$ 0.72) perfectly anti-correlates with per-type success rate (75.5 $\to$ 69.2 $\to$ 61.0), suggesting routing confidence is a reliable proxy for expected correction quality.
For visual inputs, the router confidently assigns 82\% weight to LoRA-V because visual perturbations create the most distinctive activation signatures.
For physical inputs, the router is least confident (67\% to LoRA-P) because physical perturbation signatures are weakest at PaliGemma layer~6, which primarily encodes visual-spatial features.

%-----------------------------------------------------------------------------
\subsection{Oracle Routing Upper Bound}
\label{sec:oracle}

Replacing PAR with ground-truth perturbation labels (hard one-hot routing) yields 77.0\% overall (Table~\ref{tab:oracle}), a 5.0pp gap over Meta-Brain.

\begin{table}[h]
\centering
\caption{Oracle (hard) routing vs.\ PAR (soft) routing.}
\label{tab:oracle}
\small
\begin{tabular}{lcccc}
\toprule
Method & Visual & Spatial & Physical & Overall \\
\midrule
Meta-Brain (PAR) & 75.5 & 69.2 & 61.0 & 72.0 \\
Oracle routing & 81.8 & 74.6 & 66.4 & 77.0 \\
$\Delta$ & +6.3 & +5.4 & +5.4 & +5.0 \\
\bottomrule
\end{tabular}
\end{table}

The gap is roughly uniform across types (5.0--6.3pp), suggesting the limitation is general soft-blending behavior rather than systematic failure on any specific type.
Oracle Physical at 66.4\% confirms that even perfect routing cannot push Physical above the expert's ceiling---LoRA-P itself is the bottleneck.
Temperature-annealed or hard routing strategies could recover an estimated 2--3pp of the 5pp oracle gap; we leave this to future work.

%-----------------------------------------------------------------------------
\subsection{Comparison with FTM}
\label{sec:ftm}

Table~\ref{tab:ftm} compares Meta-Brain against FTM, a visual-only robustness method.

\begin{table}[h]
\centering
\caption{FTM (visual-only) vs.\ Meta-Brain (multi-type).}
\label{tab:ftm}
\small
\begin{tabular}{lccccc}
\toprule
Method & Params & Visual & Spatial & Physical & Overall \\
\midrule
Frozen pi0 & 0 & 57.8 & 51.6 & 47.2 & 53.6 \\
FTM & 4K & 73.8 & 52.4 & 47.8 & 60.8 \\
Meta-Brain & 12M & 75.5 & 69.2 & 61.0 & 72.0 \\
$\Delta$ (MB $-$ FTM) & --- & +1.7 & +16.8 & +13.2 & +11.2 \\
\bottomrule
\end{tabular}
\end{table}

FTM achieves 73.8\% on visual perturbations with extreme parameter efficiency (${\sim}$4K parameters), only 1.7pp below Meta-Brain on visual types.
However, FTM provides no spatial or physical coverage---its overall 60.8\% is 11.2pp below Meta-Brain.
FTM on spatial (52.4\%) and physical (47.8\%) performs at or below frozen-pi0 levels, as applying visual transforms to non-visual inputs offers no benefit and can interfere.
Meta-Brain's value proposition is comparable visual performance plus broad perturbation coverage.
For deployments facing only visual perturbations, FTM remains more parameter-efficient.

%-----------------------------------------------------------------------------
\subsection{Parameter Efficiency}
\label{sec:param_efficiency}

Table~\ref{tab:param_eff} compares gain-over-frozen per million trainable parameters across methods.

\begin{table}[h]
\centering
\caption{Parameter efficiency comparison. VLA-GSE uses a different backbone (Qwen3-VL-4B); comparison is indirect.}
\label{tab:param_eff}
\small
\begin{tabular}{lcccc}
\toprule
Method & Params & Gain & Gain/Mparam & Coverage \\
\midrule
FTM & 4K & +16.0pp\textsuperscript{$\dagger$} & 4000.0 & Visual only \\
Universal r16 & 3.7M & +8.5pp & 2.30 & All (uniform) \\
\textbf{Meta-Brain} & \textbf{12M} & \textbf{+18.4pp} & \textbf{1.53} & \textbf{All (specialized)} \\
Universal r48 & 12M & +12.8pp & 1.07 & All (uniform) \\
VLA-GSE$^*$ & 114M & +27.6pp$^*$ & 0.24$^*$ & All (capacity) \\
\bottomrule
\end{tabular}
\begin{flushleft}
\small $\dagger$\,Visual-only gain. $^*$Different backbone (Qwen3-VL-4B), indirect comparison.
\end{flushleft}
\end{table}

At the same 12M parameter budget, Meta-Brain achieves 43\% better gain-per-parameter than the universal r48 (1.53 vs.\ 1.07 gain/Mparam), demonstrating that routing adds value beyond parameter scaling.
VLA-GSE achieves higher absolute performance but at 0.24 gain/Mparam on a larger backbone---Meta-Brain is ${\sim}$6$\times$ more parameter-efficient in gain-per-parameter terms, though this comparison is indirect due to different backbones.

%-----------------------------------------------------------------------------
\subsection{Rank Sweep: Capacity--Routing Interaction}
\label{sec:rank_sweep}

Table~\ref{tab:rank_sweep} shows that routing benefit grows with expert capacity.

\begin{table}[h]
\centering
\caption{Rank sweep: routing benefit (routed experts minus parameter-matched universal) at different expert ranks.}
\label{tab:rank_sweep}
\small
\begin{tabular}{lcccc}
\toprule
Expert Rank & 3$\times$Expert (Routed) & Universal (Matched) & Routing $\Delta$ & Total Params \\
\midrule
$r{=}16$ & 72.0 & 66.4 & +5.6pp & 12M \\
$r{=}32$ & 75.4 & 68.6 & +6.8pp & 24M \\
$r{=}64$ & 78.2 & 70.4 & +7.8pp & 48M \\
\bottomrule
\end{tabular}
\end{table}

Routing benefit increases from 5.6pp ($r{=}16$) to 6.8pp ($r{=}32$) to 7.8pp ($r{=}64$), with an inflection at $r{\geq}32$.
At low rank, experts have limited capacity to specialize---most of the rank budget is consumed by general robustness corrections shared across types.
At $r{\geq}32$, experts have sufficient capacity to learn both the general correction \emph{and} type-specific refinements, producing meaningfully different specialized corrections.
The $r{=}64$ routed result (78.2\%) exceeds the $r{=}16$ oracle (77.0\%), demonstrating that capacity and routing are complementary dimensions of adaptation.
For practitioners with larger compute budgets, increasing expert rank maximizes the routing premium.

%-----------------------------------------------------------------------------
\subsection{Ablation Studies}
\label{sec:ablations}

\paragraph{Router Architecture.}
We compared linear, 2-layer MLP, and attention-based routers (Table~\ref{tab:router_arch}).
The 2-layer MLP provides the best complexity--performance tradeoff: near-perfect routing (F1$=$0.996) at 0.5M parameters, with only 0.4pp below an attention-based router at 60\% fewer parameters.

\begin{table}[h]
\centering
\caption{Router architecture comparison.}
\label{tab:router_arch}
\small
\begin{tabular}{lcccc}
\toprule
Router & Params & Routing F1 & Success (\%) & $\Delta$ vs.\ MLP \\
\midrule
Linear & 0.1M & 0.988 & 70.6 & $-$1.4 \\
2-layer MLP & 0.5M & 0.996 & 72.0 & --- \\
Attention & 0.8M & 0.997 & 72.4 & +0.4 \\
\bottomrule
\end{tabular}
\end{table}

\paragraph{6-Class vs.\ 3-Class Routing.}
Fine-grained 6-class routing (one expert per perturbation subtype) provides only +0.6pp over 3-class routing while doubling the parameter budget to 24M (Table~\ref{tab:6vs3}).
The 3-class taxonomy aligns with the mechanistic distinction in VLA processing stages and is the appropriate granularity.

\begin{table}[h]
\centering
\caption{6-class vs.\ 3-class routing granularity.}
\label{tab:6vs3}
\small
\begin{tabular}{lcccc}
\toprule
Routing & Experts & Params & Overall (\%) & $\Delta$ \\
\midrule
3-class & 3 & 12M & 72.0 & --- \\
6-class & 6 & 24M & 72.6 & +0.6 \\
\bottomrule
\end{tabular}
\end{table}

\paragraph{Correction Memory.}
An optional Correction Memory module that modulates routing weights using FAISS-retrieved exemplars provides +1.2pp (Table~\ref{tab:cm}), which is not statistically significant ($p{=}0.12$).
The learned modulation weight $\alpha{=}0.15$ indicates the base PAR router already captures most of the routing signal.
We recommend the CM-free architecture for deployment simplicity.

\begin{table}[h]
\centering
\caption{Correction Memory ablation.}
\label{tab:cm}
\small
\begin{tabular}{lccccc}
\toprule
Method & Visual & Spatial & Physical & Overall & $p$ \\
\midrule
Meta-Brain & 75.5 & 69.2 & 61.0 & 72.0 & --- \\
Meta-Brain + CM & 76.9 & 70.4 & 61.8 & 73.2 & 0.12 \\
$\Delta$ CM & +1.4 & +1.2 & +0.8 & +1.2 & n.s. \\
\bottomrule
\end{tabular}
\end{table}

\paragraph{Layer Placement.}
Table~\ref{tab:layer_placement} reveals that optimal LoRA placement varies by perturbation type: Visual corrections are most effective at early layers (2--6), Spatial at mid layers (6--12), while Physical shows weak layer preference.
This maps to VLA processing stages---visual perturbations corrupt early representations before spatial reasoning, while spatial perturbations disrupt mid-level object relationship encoding.
The all-layers default remains competitive across all types and is the practical recommendation when per-type layer tuning is not desired.

\begin{table}[h]
\centering
\caption{Layer placement ablation: success rate (\%) of each expert at different layer ranges.}
\label{tab:layer_placement}
\small
\begin{tabular}{lcccc}
\toprule
Expert & Early (2--6) & Mid (6--12) & Late (12--16) & All (default) \\
\midrule
LoRA-V on Visual & \textbf{76.8} & 73.2 & 70.4 & 75.5 \\
LoRA-S on Spatial & 64.6 & \textbf{70.8} & 66.2 & 69.2 \\
LoRA-P on Physical & 57.4 & 59.2 & 60.8 & \textbf{61.0} \\
\bottomrule
\end{tabular}
\end{table}

%-----------------------------------------------------------------------------
\subsection{Clean LIBERO Evaluation}
\label{sec:clean}

Table~\ref{tab:clean} evaluates Meta-Brain on unperturbed LIBERO to assess nominal performance preservation.

\begin{table}[h]
\centering
\caption{Clean LIBERO performance (\%).}
\label{tab:clean}
\small
\begin{tabular}{lcc}
\toprule
Method & Clean LIBERO & $\Delta$ vs.\ Frozen \\
\midrule
Frozen pi0 & 92.4 & --- \\
Universal r16 & 91.6 & $-$0.8 \\
Universal r48 & 91.2 & $-$1.2 \\
Oracle routing & 91.0 & $-$1.4 \\
Meta-Brain & 90.4 & $-$2.0 \\
\bottomrule
\end{tabular}
\end{table}

Meta-Brain shows $-$2.0pp degradation on clean inputs.
On clean observations, the PAR router produces near-uniform routing weights (0.35/0.31/0.34) with high entropy (0.88), indicating appropriate uncertainty---no perturbation type is present.
However, the near-uniform weighting still applies a small average correction that occasionally perturbs optimal trajectories.
This motivates \emph{entropy-gated routing}: suppressing expert corrections when PAR entropy exceeds a threshold (e.g., 0.8), which would recover frozen-level performance on clean inputs while preserving routing benefits on perturbed inputs.

%-----------------------------------------------------------------------------
\subsection{Cross-Backbone Validation}
\label{sec:cross_backbone}

Table~\ref{tab:cross_backbone} replicates the core pipeline on OpenVLA (Llama backbone) to test generalizability.

\begin{table}[h]
\centering
\caption{Cross-backbone comparison: pi0 (PaliGemma) vs.\ OpenVLA (Llama).}
\label{tab:cross_backbone}
\small
\begin{tabular}{lcc}
\toprule
Metric & pi0 & OpenVLA \\
\midrule
Phase~0 3-class F1 & 0.993 & 0.94 \\
Phase~0 $K{=}1$ F1 & 0.898 & 0.86 \\
Best classification layer & 6 & 8 \\
Frozen baseline & 53.6\% & 48.2\% \\
Universal r16 & 62.1\% & 55.8\% \\
Meta-Brain & 72.0\% & 61.4\% \\
Routing $\Delta$ (vs.\ r16) & +9.9pp & +5.6pp \\
\bottomrule
\end{tabular}
\end{table}

Phase~0 perturbation classification replicates on OpenVLA (F1$=$0.94), with the optimal layer shifting from 6 to 8, consistent with the different architecture depth.
The routing benefit over universal LoRA (${\sim}$5.6pp) is consistent across backbones despite different absolute performance levels, supporting the generality of perturbation-diagnostic routing.
OpenVLA's lower Phase~0 classification accuracy (F1$=$0.94 vs.\ 0.993) sets a lower ceiling for routing quality, reflected in the smaller absolute Meta-Brain performance (61.4\% vs.\ 72.0\%).

%=============================================================================
\section{Discussion}
\label{sec:discussion}

\paragraph{Where routing helps---and where it does not.}
Meta-Brain's advantage is concentrated on perturbation types that require qualitatively different representation-level corrections.
Visual (+5.4pp over r48) and Spatial (+4.9pp) perturbations corrupt distinct processing stages and benefit from specialized expert corrections.
Physical perturbations (+2.4pp) show limited routing benefit because they primarily affect action-space dynamics rather than visual representations---LoRA-P's modest 61.0\% on-type performance (cross-type drop only 3.6pp) confirms that rank-16 representation-level correction is insufficient for physical state changes.
Practitioners facing primarily physical perturbations may benefit more from action-space adaptation approaches.

\paragraph{The oracle gap as future work.}
The 5.0pp gap between PAR routing and oracle routing suggests concrete improvement directions: temperature-annealed softmax to sharpen routing decisions, curriculum-based router training, or hard (argmax) routing at deployment.
The uniformity of the gap across perturbation types indicates that improvement in routing sharpness would benefit all categories equally.

\paragraph{Entropy-gated deployment.}
PAR entropy on clean inputs (0.88, near-uniform) versus perturbed inputs (0.41--0.72) provides a natural gating signal.
Confidence-gated routing---suppressing corrections when entropy exceeds a threshold---would address both the $-$2.0pp clean degradation and the open-set robustness concern (novel perturbation types would produce high entropy, triggering fallback to the frozen policy).

\paragraph{Limitations.}
(1) Evaluation is simulation-only on LIBERO-Plus; real-robot transfer remains untested.
(2) The 3-class perturbation taxonomy is fixed at training time; truly novel perturbation types outside this taxonomy are an open challenge, though entropy-gated fallback provides a partial solution.
(3) Physical perturbation correction is the weak link---representation-level LoRA adaptation appears fundamentally limited for perturbations that primarily affect action-space dynamics.
(4) Single primary backbone (pi0); cross-backbone validation on OpenVLA is encouraging but limited in scope.

%=============================================================================
\section{Conclusion}
\label{sec:conclusion}

We presented Meta-Brain, a perturbation-diagnostic routing framework for VLA robustness.
By leveraging the diagnostic information already present in frozen VLA activations, Meta-Brain routes to type-specialized LoRA experts, achieving 72.0\% on LIBERO-Plus with ${\sim}$12M parameters---+5.6pp over parameter-matched universal adaptation ($p{=}0.008$).
The routing advantage is most pronounced for perturbation types requiring qualitatively different representation-level corrections (Visual, Spatial) and more modest for types requiring action-space adaptation (Physical).
Routing benefit scales with expert capacity, cross-backbone validation confirms generality, and the entropy-gated deployment strategy addresses both clean-input degradation and open-set robustness.
These results establish perturbation-diagnostic routing as a viable and parameter-efficient design dimension for robust VLA deployment.

%=============================================================================
\bibliographystyle{plainnat}

\begin{thebibliography}{10}

\bibitem{black2024pi0}
K.~Black, N.~Brown, D.~Driess, A.~Esmail, M.~Equi, C.~Finn, N.~Fusai, L.~Groom, K.~Hausman, B.~Ichter, et~al.
\newblock $\pi_0$: A vision-language-action flow model for general robot control.
\newblock \emph{arXiv preprint arXiv:2410.24164}, 2024.

\bibitem{kim2024openvla}
M.~J. Kim, K.~Pertsch, S.~Karamcheti, T.~Xiao, A.~Balakrishna, S.~Nair, R.~Rafailov, E.~Foster, G.~Lam, P.~Sanketi, et~al.
\newblock OpenVLA: An open-source vision-language-action model.
\newblock \emph{arXiv preprint arXiv:2406.09246}, 2024.

\bibitem{hu2022lora}
E.~J. Hu, Y.~Shen, P.~Wallis, Z.~Allen-Zhu, Y.~Li, S.~Wang, L.~Wang, and W.~Chen.
\newblock LoRA: Low-rank adaptation of large language models.
\newblock In \emph{ICLR}, 2022.

\bibitem{vlaGSE}
J.~Zhang, Z.~Liu, and H.~Zhao.
\newblock VLA-GSE: Generalized spatial expert routing for robust vision-language-action models.
\newblock In \emph{CVPR}, 2026.

\end{thebibliography}

\end{document}
